\documentclass[10pt, a4paper]{article}
\usepackage[T2A]{fontenc}
\usepackage[russian]{babel}
\usepackage{geometry}
\usepackage{titlesec}
\usepackage{tikz}
\usepackage{amsmath}
\usepackage{amssymb}
\usepackage{array}
\usepackage{tcolorbox}
\usepackage{stackengine}
\usepackage{verbatim}
\usepackage{fancyvrb}
\usepackage{listings}
\usepackage{tocloft}
\usepackage[colorlinks=true, linkcolor=black, urlcolor=blue, citecolor=blue]{hyperref}
\usepackage{biblatex}

\geometry{top=2cm, bottom=2cm, left=2cm, right=2cm}

\setlength{\parindent}{0pt}

\definecolor{codegreen}{rgb}{0,0.6,0}
\definecolor{codegray}{rgb}{0.5,0.5,0.5}
\definecolor{codepurple}{rgb}{0.58,0,0.82}
\definecolor{backcolour}{rgb}{0.95,0.95,0.92}
\lstdefinestyle{mystyle}{
    backgroundcolor=\color{backcolour},
    commentstyle=\color{codegreen},
    keywordstyle=\color{magenta},
    numberstyle=\tiny\color{codegray},
    basicstyle=\ttfamily\footnotesize,
    breakatwhitespace=false,
    breaklines=true,
    captionpos=b,
    keepspaces=true,
    numbers=left,
    numbersep=5pt,
    showspaces=false,
    showstringspaces=false,
    showtabs=false,
    tabsize=4
}
\lstset{style=mystyle}
\lstset{extendedchars=\true}

\titleformat{\section}[block]{\centering\Large\bfseries}{}{0pt}{}
\titleformat{\subsection}[block]{\flushleft\normalsize\bfseries}{}{0pt}{}
\setcounter{secnumdepth}{-1}

\renewcommand{\cftsecleader}{\cftdotfill{\cftdotsep}}
\renewcommand{\cftsubsecleader}{\cftdotfill{\cftdotsep}}
\renewcommand{\cftsubsubsecleader}{\cftdotfill{\cftdotsep}}

\newcommand{\algquote}[1]{<<#1>>}
\newcommand{\algindex}[1]{[#1]}
\newcommand{\algline}[1]{\underline{#1}}
\newcommand{\algplus}{+}
\newcommand{\algleq}{$\leq$}

\newcommand{\specframeimpl}[1]
{
\begin{array}{l}
    \begin{tcolorbox}[
        width=\textwidth,
        colback=white,
        colframe=black,
        sharp corners,
        boxrule=1pt
    ]
    #1
    \end{tcolorbox}
\end{array}
}

\newcommand{\specframe}[1]
{
\[
    \specframeimpl{#1}
\]
}

\newcommand{\specframeif}[1]
{
\[
\hspace{-2.85mm}
\left\{
    \specframeimpl{#1}
\right.
\]
}

\newcommand{\specframeloop}[1]
{
\[
\hspace{-3.12mm}
\left\{
    \specframeimpl{#1}
\right\}^{\!\bigstar}
\]
}

\newcommand{\formaltaskimpl}[2]
{
\underline{#1}
\begin{enumerate}
    #2
\end{enumerate}
}

\newcommand{\given}[1]{\formaltaskimpl{Дано:}{#1}}
\newcommand{\result}[1]{\formaltaskimpl{Результат:}{#1}}
\newcommand{\prerequisites}[1]{\formaltaskimpl{При}{#1}}
\newcommand{\binding}[1]{\formaltaskimpl{Связь}{#1}}

\newcommand{\mathsolution}[1]
{
\begin{enumerate}
    #1
\end{enumerate}
}

\newcommand{\mathsolutiontask}[1]
{
\item
$
\begin{array}{l}
    #1
\end{array}
$
}

\newcommand{\mathsolutionloop}[1]
{
\begin{cases}
    #1
\end{cases}
}

\newcommand{\pseudocode}[1]
{
\begin{Verbatim}[commandchars=+\[\]]
    #1
\end{Verbatim}
}

\title{
    \textbf{Федеральное государственное автономное образовательное учреждение высшего образования}\\
    \textbf{«Национальный исследовательский университет}\\
    \textbf{"Высшая школа экономики"»}

    \vspace{1cm}

    \textbf{Московский институт электроники и математики им. А.Н. Тихонова НИУ ВШЭ}

    \textbf{Департамент компьютерной инженерии}

    \vspace{1.5cm}

    \textbf{Курс: Алгоритмизация и программирование}

    \vspace{1cm}

    \textbf{ОТЧЕТ}

    \textbf{по лабораторной работе №\labnumber}
}

\newcommand{\labnumber}{номер лабы}
\newcommand{\tasks}{(задачи)}

\author{
    \fontsize{10}{10}
    \hspace{-1cm}
    \begin{tabular}{cc}
        \begin{tabular}{|l|c|c|}
            \hline
            \textbf{Раздел} & \textbf{Mах оценка} & \textbf{Итог. оценка} \\
            \hline
            Постановка & $0,5$ & \\
            \hline
            Метод & $1$ & \\
            \hline
            Спецификация & $0,5$ & \\
            \hline
            Алгоритм & $1,5$ & \\
            \hline
            Работа программы & $1$ & \\
            \hline
            Листинг & $0,5$ & \\
            \hline
            Тесты & $1$ & \\
            \hline
            Вопросы & $2$ & \\
            \hline
            Доп. задание & 2 & \\
            \hline
        \end{tabular}
        &
        \begin{tabular}{l}
            Студент: Солодилов Михаил Анатольевич \\
            Группа: БИВ253 \\
            Вариант: 232 \tasks\\
            Руководитель:\\
            Оценка:\\
            Дата сдачи \\
        \end{tabular}
    \end{tabular}
}

\date{
    \vfill
    \textbf{МОСКВА 2025}
}

\addto\captionsrussian{\renewcommand\contentsname{Оглавление}}


\renewcommand{\labnumber}{2}
\renewcommand{\tasks}{(12, 3)}

\begin{document}

\maketitle
\thispagestyle{empty}
\newpage
\pagenumbering{arabic}
\tableofcontents

\newpage
\section{Задание}

\begin{enumerate}
    \item Даны целочисленная матрица X[1:n, 1:m] и целочисленный
        массив Z[1:k]. Обнулить элементы матрицы X, которых нет в
        массиве Z и запомнить обнуленные элементы.
    \item Дан массив целых положительных чисел. Сформировать новый
        массив, содержащий произведения цифр каждого элемента
        исходного массива.
\end{enumerate}

\newpage
\section{Постановка задачи}

\given{
    \item
        n, m, k --- цел. \\
        x[1:n, 1:m] --- цел. \\
        z[1:k] --- цел.
    \item
        n - цел.
        a[1:n] --- цел.
}
\result{
    \item
        t --- цел., количество обнулённых элементов \\
        e[1:t] --- обнулённые элементы или сообщение <<Ноль
        обнулений>>
    \item
        p[1:n] --- цел.
}
\prerequisites{
    \item
        $ 1 \leq n, m, k \leq LMAX$ \\
    \item
        $ 1 \leq n \leq LMAX$
}
\binding{
    \item
        $\forall i \in [1, k]: \exists (i, j) \in \{n \times m\}:
        x[i, j] = z[i] \implies x'[i, j] = 0, \exists k \in [1, t]:
        e[t] = z[i]$, где $x'$ --- массив после преобразования
    \item
        $\forall i \in [1, n] \implies p[i] = productdigits(a[i])$
}

\newpage
\section{Метод решения задачи}

\mathsolution{
    \mathsolutiontask{
        t = 0 \\
        \mathsolutionloop{
            \text{для } i = \overline{1, n} \\
            \mathsolutionloop{
                \text{для } j = \overline{1, n} \\
                found = 1 \\
                \mathsolutionloop{
                    \text{пока } found \leq k \text{ и } x[i, j]
                    \neq z[found] \\
                    found = found + 1
                } \\
                x[i, j] = 0 \text{, если } found \leq k \\
                t = t + 1 \text{, если } found \leq k \\
                e[t] = z[found] \text{, если } found \leq k
            }
        }
    }
    \mathsolutiontask{
        \mathsolutionloop{
            \text{для } i = \overline{1, n} \\
            t_{prod} = 1 \\
            t_{val} = a[i] \\
            \mathsolutionloop{
                \text{пока } t_{val} > 0 \\
                t_{prod} = t_{prod} \cdot (t_{val} \bmod 10) \\
                t_{val} = \lfloor t_{val} / 10 \rfloor \\
            } \\
            p[i] = t_{prod}
        }
    }
}

\newpage
\section{Внешняя спецификация}

\specframe{
    Лабораторная работа №\labnumber
}

\specframe{
    Задание №1
}

\specframeloop{
    Введите n, m, k от 1 до LMAX: \\
    $<n>$, $<m>$, $<k>$
}
До $1 \leq n \leq LMAX$ и $1 \leq m \leq LMAX$ и $1 \leq k \leq LMAX$

\specframe{
    Введите матрицу x: \\
    $\ll x[1, 1] \gg$, $\ll x[1, 2] \gg$, ..., $\ll x[2, 1] \gg$, $\ll x[2, 2] \gg$, ..., $\ll x[n, m] \gg$,
}

\specframe{
    Задание №1
}

\newpage
\section{Описание алгоритма на псевдокоде}

% +algquote[]
% +algindex[]
% +algline[]
% +algplus
% +algleq

\begin{Verbatim}[commandchars=+\[\]]
\end{Verbatim}

\newpage
\section{Листинг программы}

% \lstinputlisting[language=C]{main.c}

\newpage
\section{Распечатка тестов к программе и результатов}

\end{document}
